\documentclass[10pt,a4paper]{amsart}
\usepackage[margin=19mm]{geometry}
\usepackage{amsmath,amssymb,mathtools}
\usepackage[scr=rsfs,cal=euler]{mathalfa}
\usepackage{xfrac}
\usepackage[unicode,hidelinks,bookmarks=false]{hyperref}
\newcommand{\ie}{i.e.\ }
\newcommand{\cf}{cf.\ }
\newcommand{\resp}{respectively\ }
\newcommand{\Subsection}{\subsection}
\providecommand{\nobreakdash}{\nobreak\hbox{-}}
\setlength{\emergencystretch}{2em}
\multlinegap0pt
\usepackage{tikz}
\newtheorem{theorem}{Theorem}
\title{Support Tokens, Stability Margins, and \\ a New Foundation for Robust LLMs \vspace{1em}}
\author{Deepak Agarwal, Dhyey Dharmendrakumar Mavani, Suyash Gupta, Karthik Sethuraman, Tejas Dharamsi}
\date{Source: arXiv:2602.22271v3 (CC BY 4.0)}
\begin{document}
\maketitle

\noindent\emph{Source and licence.} Support Tokens, Stability Margins, and \\ a New Foundation for Robust LLMs \vspace{1em}. Version arXiv:2602.22271v3 (CC BY 4.0), distributed by arXiv under the Creative Commons Attribution 4.0 International licence, http://creativecommons.org/licenses/by/4.0/, as recorded in the arXiv licence metadata for this exact version and shown on the abstract page https://arxiv.org/abs/2602.22271v3. Full licence notice: \texttt{licenses/arxiv-2602.22271-CC-BY-4.0.txt}.\par\smallskip

\noindent\textit{Editorial scope.} Counted unit: the article's theorem theorem (source0.tex lines 2264--2271). The article's complete original proof of it is delivered with it (source0.tex lines 2273--2313, one \texttt{proof} environment of 41 lines). No mathematics is written, repaired or re-derived: the statement and the proof are the article's own lines, in the article's own notation, and no retained line is substituted or re-flowed.  Retained uncounted as the source-local context the proof needs: lines 706--714.  Nothing was omitted from any retained span, so the package carries no omission record.  No retained line contains a citation, so the wrapper carries no bibliography and no bibliography entry.  No retained line contains a figure environment, an image-inclusion command or drawing code, so no image is delivered and the package declares no asset dependency.  Editorial formatting only, in addition: an AMS \texttt{amsart} preamble that restates the article's own declarations and macros used by the retained lines, namely the environments \texttt{theorem} on separate counters, together with the standard packages those lines need.  The retained environments print in this document as Theorem 1 in order of appearance, whereas the article prints the same items with the numbering of its own full document.  The complete native source of the article as distributed in the arXiv e-print for this exact version is bundled unchanged as \texttt{sources/195321/source0.tex}.
\begin{theorem}[Kolmogorov consistency]
\label{thm:kolmogorov_consistency}
Assume strict causal masking and Jacobian non-degeneracy so that $p_\sigma(x_{1:n})$ exists for each $n$. Then for all $n\ge 2$ and all $y_{1:n-1}\in V^{n-1}$,
\begin{equation}
\sum_{y_n\in V} P(y_{1:n}) \;=\; P(y_{1:n-1}),
\label{eq:kolmogorov_consistency}
\end{equation}
where the right-hand side is the marginal likelihood obtained from the \emph{$(n-1)$-dimensional} model.
\end{theorem}

\begin{theorem}
Under the assumptions of Theorem~\ref{thm:kolmogorov_consistency}, there exists a unique probability
measure $P$ on $(V^\infty,\mathcal F)$ whose finite-dimensional marginals are $\{P_n\}_{n\ge 1}$:
for each $n$ and each $y_{1:n}\in V^n$,
\[
P(\{Y_{1:n}=y_{1:n}\}) = P_n(y_{1:n}).
\]
\end{theorem}

\begin{proof}
We verify the hypotheses of the Kolmogorov extension theorem for the projective family $\{P_n\}_{n\ge 1}$.

\textbf{Step 1.}
For each $n$ and each $y_{1:n}\in V^n$, define the cylinder set
\[
C(y_{1:n}) \;\triangleq\; \{ \omega\in V^\infty : \omega_1=y_1,\dots,\omega_n=y_n\}.
\]
Define $\nu$ on cylinders by $\nu(C(y_{1:n})) \triangleq P_n(y_{1:n})$.

\textbf{Step 2.}
If a cylinder is described at two different lengths, consistency ensures the same value.
Indeed, for any fixed $y_{1:n}$,
\[
\nu(C(y_{1:n}))
=
P_n(y_{1:n})
=
\sum_{y_{n+1}\in V} P_{n+1}(y_{1:n},y_{n+1})
=
\sum_{y_{n+1}} \nu(C(y_{1:n},y_{n+1})),
\]
so the mass assigned to a cylinder equals the sum of masses of its disjoint refinements.

\textbf{Step 3.}
The set of finite disjoint unions of cylinders forms an algebra that generates $\mathcal F$.
Using Step 2 (refinement) and the fact that each $P_n$ is a probability measure on the finite set $V^n$,
$\nu$ is finitely additive on this algebra.

\textbf{Step 4.}
Because $V$ is finite, the cylinder algebra is countable, and $\nu$ is automatically $\sigma$-additive
on it (one can check directly by refining to a common length and using $\sigma$-additivity on $V^n$).
Therefore $\nu$ is a pre-measure on a generating algebra of $\mathcal F$.
By Carath\'eodory's extension theorem, $\nu$ extends to a measure $P$ on $(V^\infty,\mathcal F)$.

\textbf{Step 5.}
Uniqueness follows because cylinder sets generate $\mathcal F$ and any two measures agreeing on a generating
$\pi$-system agree on the generated $\sigma$-algebra.

This proves existence and uniqueness of $P$ with the prescribed marginals.
\end{proof}

\end{document}
